\SourceLeaf{208}\begin{center}
[\textsc{Bessel} \emph{an} \textsc{Gauss}. \emph{Königsberg i. Pr., 10. Februar 1829.]}
\end{center}

\{.\ .\ .\ .\ .\ .\ Ich würde sehr beklagen, wenn Sie Sich »durch das Geschrei der
Böotier« abhalten liessen, Ihre geometrischen Ansichten aus einander zu setzen.
Durch das, was \textsc{Lambert} gesagt hat, und was \textsc{Schweikart} mündlich äusserte,
ist mir klar geworden, dass unsere Geometrie unvollständig ist, und eine Correction erhalten sollte, welche hypothetisch ist und, wenn die Summe der
Winkel des ebenen Dreiecks $=180^{0}$ ist, verschwindet. Das wäre die \textls{wahre}
Geometrie, die Euklidische die \textls{praktische}, wenigstens für Figuren auf der
Erde. .\ .\ .\ .\ .\ .\}

\rule{4cm}{0.5pt}

\Needspace{16\baselineskip}
\begin{center}
\textsc{Gauss} an \textsc{Bessel}. Göttingen, 9. April 1830.
\end{center}

.\ .\ .\ .\ .\ .\ .\ Wahre Freude hat mir die Leichtigkeit gemacht, mit der Sie
in meine Ansichten über die Geometrie eingegangen sind, zumal da so wenige
offenen Sinn dafür haben. Nach meiner innigsten Überzeugung hat die Raumlehre in unserm Wissen a priori eine ganz andere Stellung, wie die reine
Grössenlehre; es geht unserer Kenntniss von jener durchaus diejenige vollständige Überzeugung von ihrer Nothwendigkeit (also auch von ihrer absoluten
Wahrheit) ab, die der letztern eigen ist; wir müssen in Demuth zugeben,
dass, wenn die Zahl \textls{bloss} unsers Geistes Product ist, der Raum auch ausser
unserm Geiste eine Realität hat, der wir a priori ihre Gesetze nicht vollständig vorschreiben können. .\ .\ .\ .\ .\ .\ .\ .

\rule{4cm}{0.5pt}

